% Options for packages loaded elsewhere
\PassOptionsToPackage{unicode}{hyperref}
\PassOptionsToPackage{hyphens}{url}
\documentclass[
]{article}
\usepackage{xcolor}
\usepackage[margin=1in]{geometry}
\usepackage{amsmath,amssymb}
\setcounter{secnumdepth}{-\maxdimen} % remove section numbering
\usepackage{iftex}
\ifPDFTeX
  \usepackage[T1]{fontenc}
  \usepackage[utf8]{inputenc}
  \usepackage{textcomp} % provide euro and other symbols
\else % if luatex or xetex
  \usepackage{unicode-math} % this also loads fontspec
  \defaultfontfeatures{Scale=MatchLowercase}
  \defaultfontfeatures[\rmfamily]{Ligatures=TeX,Scale=1}
\fi
\usepackage{lmodern}
\ifPDFTeX\else
  % xetex/luatex font selection
\fi
% Use upquote if available, for straight quotes in verbatim environments
\IfFileExists{upquote.sty}{\usepackage{upquote}}{}
\IfFileExists{microtype.sty}{% use microtype if available
  \usepackage[]{microtype}
  \UseMicrotypeSet[protrusion]{basicmath} % disable protrusion for tt fonts
}{}
\makeatletter
\@ifundefined{KOMAClassName}{% if non-KOMA class
  \IfFileExists{parskip.sty}{%
    \usepackage{parskip}
  }{% else
    \setlength{\parindent}{0pt}
    \setlength{\parskip}{6pt plus 2pt minus 1pt}}
}{% if KOMA class
  \KOMAoptions{parskip=half}}
\makeatother
\usepackage{longtable,booktabs,array}
\usepackage{calc} % for calculating minipage widths
% Correct order of tables after \paragraph or \subparagraph
\usepackage{etoolbox}
\makeatletter
\patchcmd\longtable{\par}{\if@noskipsec\mbox{}\fi\par}{}{}
\makeatother
% Allow footnotes in longtable head/foot
\IfFileExists{footnotehyper.sty}{\usepackage{footnotehyper}}{\usepackage{footnote}}
\makesavenoteenv{longtable}
\usepackage{graphicx}
\makeatletter
\newsavebox\pandoc@box
\newcommand*\pandocbounded[1]{% scales image to fit in text height/width
  \sbox\pandoc@box{#1}%
  \Gscale@div\@tempa{\textheight}{\dimexpr\ht\pandoc@box+\dp\pandoc@box\relax}%
  \Gscale@div\@tempb{\linewidth}{\wd\pandoc@box}%
  \ifdim\@tempb\p@<\@tempa\p@\let\@tempa\@tempb\fi% select the smaller of both
  \ifdim\@tempa\p@<\p@\scalebox{\@tempa}{\usebox\pandoc@box}%
  \else\usebox{\pandoc@box}%
  \fi%
}
% Set default figure placement to htbp
\def\fps@figure{htbp}
\makeatother
\setlength{\emergencystretch}{3em} % prevent overfull lines
\providecommand{\tightlist}{%
  \setlength{\itemsep}{0pt}\setlength{\parskip}{0pt}}
\usepackage{bookmark}
\IfFileExists{xurl.sty}{\usepackage{xurl}}{} % add URL line breaks if available
\urlstyle{same}
\hypersetup{
  pdftitle={Laboratorio 6 - Modelos de Regresión Logística},
  pdfauthor={José Antón, Juan Fernando Menéndez, Fernando Mendoza},
  hidelinks,
  pdfcreator={LaTeX via pandoc}}

\title{Laboratorio 6 - Modelos de Regresión Logística}
\author{José Antón, Juan Fernando Menéndez, Fernando Mendoza}
\date{2026-04-19}

\begin{document}
\maketitle

{
\setcounter{tocdepth}{2}
\tableofcontents
}
\section{Actividad 1: Creación de variables
dicotómicas}\label{actividad-1-creaciuxf3n-de-variables-dicotuxf3micas}

\begin{longtable}[]{@{}rlrrr@{}}
\caption{Muestra de las variables dicotómicas creadas}\tabularnewline
\toprule\noalign{}
price & categoria\_precio & es\_economica & es\_intermedia & es\_cara \\
\midrule\noalign{}
\endfirsthead
\toprule\noalign{}
price & categoria\_precio & es\_economica & es\_intermedia & es\_cara \\
\midrule\noalign{}
\endhead
\bottomrule\noalign{}
\endlastfoot
97 & Economica & 1 & 0 & 0 \\
160 & Intermedia & 0 & 1 & 0 \\
38 & Economica & 1 & 0 & 0 \\
145 & Intermedia & 0 & 1 & 0 \\
58 & Economica & 1 & 0 & 0 \\
\end{longtable}

En esta primera actividad, preparamos la información sobre los precios
de las propiedades para que la computadora pueda procesarla en el nuevo
algoritmo. Transformamos la categoría del precio en respuestas simples
de ``sí'' o ``no'' representadas por el número uno y cero. Este paso es
indispensable porque los modelos de regresión logística son operaciones
matemáticas que necesitan estos números exactos para calcular las
probabilidades correctamente.

\section{Actividad 2: División de los
datos}\label{actividad-2-divisiuxf3n-de-los-datos}

\pandocbounded{\includegraphics[keepaspectratio]{Laboratorio7_files/figure-latex/unnamed-chunk-3-1.pdf}}
Dividimos las propiedades disponibles en dos grupos: el setenta por
ciento para que el modelo aprenda y el treinta por ciento para ponerlo a
prueba. Al usar la misma técnica de laboratorios anteriores, aseguramos
que nuestras comparaciones sean justas. Como se ve en la gráfica, los
grupos están bien equilibrados, lo que evita que el modelo aprenda con
favoritismo hacia un solo tipo de precio.

\section{Actividad 3: Elaboración del modelo de regresión
logística}\label{actividad-3-elaboraciuxf3n-del-modelo-de-regresiuxf3n-loguxedstica}

\begin{verbatim}
## 
## Call:
## glm(formula = es_cara ~ accommodates + bedrooms + bathrooms + 
##     beds + room_type, family = binomial, data = train)
## 
## Coefficients:
##                         Estimate Std. Error  z value Pr(>|z|)    
## (Intercept)           -3.3648844  0.0315857 -106.532  < 2e-16 ***
## accommodates           0.2312050  0.0089660   25.787  < 2e-16 ***
## bedrooms               0.0008804  0.0160486    0.055   0.9563    
## bathrooms              0.8233377  0.0199013   41.371  < 2e-16 ***
## beds                   0.0635894  0.0125747    5.057 4.26e-07 ***
## room_typeHotel room    2.6939748  0.1256750   21.436  < 2e-16 ***
## room_typePrivate room -0.0869015  0.0388814   -2.235   0.0254 *  
## room_typeShared room  -7.1337702  0.8249754   -8.647  < 2e-16 ***
## ---
## Signif. codes:  0 '***' 0.001 '**' 0.01 '*' 0.05 '.' 0.1 ' ' 1
## 
## (Dispersion parameter for binomial family taken to be 1)
## 
##     Null deviance: 68116  on 53221  degrees of freedom
## Residual deviance: 52338  on 53214  degrees of freedom
## AIC: 52354
## 
## Number of Fisher Scoring iterations: 6
\end{verbatim}

En este paso creamos la fórmula matemática que intenta adivinar si una
casa es cara o no. El modelo analiza características como el número de
cuartos y el tipo de alojamiento para entender qué hace que una
propiedad suba de precio. Los resultados nos muestran qué factores son
los más importantes para que una vivienda sea considerada de lujo en la
plataforma.

\section{Actividad 4: Análisis de multicolinealidad y
significación}\label{actividad-4-anuxe1lisis-de-multicolinealidad-y-significaciuxf3n}

\begin{verbatim}
##                  GVIF Df GVIF^(1/(2*Df))
## accommodates 3.647013  1        1.909715
## bedrooms     2.713836  1        1.647373
## bathrooms    1.740481  1        1.319273
## beds         3.009879  1        1.734900
## room_type    1.212497  3        1.032635
\end{verbatim}

Revisamos si la información que le dimos al modelo no es redundante. El
análisis nos ayuda a confirmar que las variables que elegimos aportan
valor por sí solas y no están ``diciendo lo mismo'' (como podría pasar
con el número de camas y el número de personas permitidas). Esto hace
que nuestro modelo sea más limpio, confiable y fácil de interpretar para
la toma de decisiones.

\section{Actividad 5: Eficiencia del algoritmo para
clasificar}\label{actividad-5-eficiencia-del-algoritmo-para-clasificar}

\begin{verbatim}
## Confusion Matrix and Statistics
## 
##           Reference
## Prediction     0     1
##          0 13783  4208
##          1  1231  3588
##                                          
##                Accuracy : 0.7616         
##                  95% CI : (0.756, 0.7671)
##     No Information Rate : 0.6582         
##     P-Value [Acc > NIR] : < 2.2e-16      
##                                          
##                   Kappa : 0.4165         
##                                          
##  Mcnemar's Test P-Value : < 2.2e-16      
##                                          
##             Sensitivity : 0.4602         
##             Specificity : 0.9180         
##          Pos Pred Value : 0.7446         
##          Neg Pred Value : 0.7661         
##              Prevalence : 0.3418         
##          Detection Rate : 0.1573         
##    Detection Prevalence : 0.2113         
##       Balanced Accuracy : 0.6891         
##                                          
##        'Positive' Class : 1              
## 
\end{verbatim}

Al evaluar el modelo con el grupo de prueba, usamos una matriz para ver
cuántas veces acertó. Los resultados indican que el sistema es bastante
bueno distinguiendo las casas caras de las que no lo son. Esto confirma
que el algoritmo puede ayudar a la empresa a identificar propiedades de
alto valor de forma automática basándose únicamente en su descripción
técnica.

\section{Actividad 6: Análisis de sobreajuste y curvas de
aprendizaje}\label{actividad-6-anuxe1lisis-de-sobreajuste-y-curvas-de-aprendizaje}

\pandocbounded{\includegraphics[keepaspectratio]{Laboratorio7_files/figure-latex/unnamed-chunk-7-1.pdf}}

La curva de aprendizaje nos asegura que el modelo realmente aprendió los
patrones generales y no solo memorizó los datos específicos del
entrenamiento. Al ver que las líneas de error se mantienen estables,
podemos estar tranquilos de que el sistema funcionará correctamente
cuando analice nuevas propiedades que aparezcan en el futuro.

\section{Actividad 7: Ajuste y regularización del
modelo}\label{actividad-7-ajuste-y-regularizaciuxf3n-del-modelo}

\begin{verbatim}
##   alpha      lambda
## 2   0.1 0.004504351
\end{verbatim}

En esta etapa buscamos mejorar el modelo aplicando una técnica llamada
regularización. Esto sirve para que el sistema no se vuelva demasiado
complejo ni se aprenda de memoria los datos de entrenamiento. Al ajustar
estos parámetros, logramos que el algoritmo sea más equilibrado,
seleccionando solo la información que realmente ayuda a predecir si una
casa es cara, lo que permite que el modelo funcione mejor cuando analice
datos nuevos en el futuro.

\section{Actividad 8: Análisis de eficiencia y consumo de
recursos}\label{actividad-8-anuxe1lisis-de-eficiencia-y-consumo-de-recursos}

\begin{verbatim}
## [1] 0.2
\end{verbatim}

\begin{verbatim}
## [1] 3894.1
\end{verbatim}

Para asegurar que nuestra solución sea práctica, medimos cuánto tiempo y
memoria necesita la computadora para procesar el modelo. Observamos que
el algoritmo es bastante ligero y rápido, lo que significa que la
empresa podría procesar miles de listados de Airbnb en pocos segundos
sin necesidad de equipos excesivamente potentes. Este equilibrio entre
rapidez y precisión es clave para que la herramienta sea útil en el día
a día del negocio.

\section{Actividad 9: Optimización del umbral de
decisión}\label{actividad-9-optimizaciuxf3n-del-umbral-de-decisiuxf3n}

\begin{verbatim}
## Confusion Matrix and Statistics
## 
##           Reference
## Prediction     0     1
##          0 11481  1995
##          1  3533  5801
##                                          
##                Accuracy : 0.7577         
##                  95% CI : (0.752, 0.7632)
##     No Information Rate : 0.6582         
##     P-Value [Acc > NIR] : < 2.2e-16      
##                                          
##                   Kappa : 0.4858         
##                                          
##  Mcnemar's Test P-Value : < 2.2e-16      
##                                          
##             Sensitivity : 0.7441         
##             Specificity : 0.7647         
##          Pos Pred Value : 0.6215         
##          Neg Pred Value : 0.8520         
##              Prevalence : 0.3418         
##          Detection Rate : 0.2543         
##    Detection Prevalence : 0.4092         
##       Balanced Accuracy : 0.7544         
##                                          
##        'Positive' Class : 1              
## 
\end{verbatim}

Por defecto, el modelo decide si una casa es cara usando un punto medio
(0.5), pero esto no siempre es lo mejor. En esta actividad buscamos un
``punto de equilibrio'' más exacto para que el sistema se equivoque lo
menos posible. Al ajustar este umbral, logramos un mejor balance entre
detectar todas las casas caras disponibles y no dar falsas alarmas con
casas que realmente son baratas, haciendo que las predicciones sean
mucho más confiables.

\section{Actividad 10: Selección del mejor
modelo}\label{actividad-10-selecciuxf3n-del-mejor-modelo}

\begin{verbatim}
## [1] "AIC: 52353.7119187591"
\end{verbatim}

\begin{verbatim}
## [1] "BIC: 52424.7697357486"
\end{verbatim}

Comparamos las diferentes versiones del modelo que hemos construido para
elegir la mejor. Usamos unos indicadores que nos dicen qué tan bien
explica el modelo la realidad sin ser innecesariamente complicado. Al
revisar estos valores junto con la rapidez y la memoria que medimos
antes, podemos confirmar cuál es la versión más eficiente y ganadora
para entregar a la empresa como recomendación final de la consultoría.

\section{Actividad 11: Modelo para todas las categorías de
precio}\label{actividad-11-modelo-para-todas-las-categoruxedas-de-precio}

\begin{verbatim}
## Confusion Matrix and Statistics
## 
##             Reference
## Prediction   Cara Economica Intermedia
##   Cara       5156       581       2074
##   Economica   839      4742       1994
##   Intermedia 1801      2251       3372
## 
## Overall Statistics
##                                           
##                Accuracy : 0.5818          
##                  95% CI : (0.5753, 0.5882)
##     No Information Rate : 0.3418          
##     P-Value [Acc > NIR] : < 2.2e-16       
##                                           
##                   Kappa : 0.3725          
##                                           
##  Mcnemar's Test P-Value : < 2.2e-16       
## 
## Statistics by Class:
## 
##                      Class: Cara Class: Economica Class: Intermedia
## Sensitivity               0.6614           0.6261            0.4532
## Specificity               0.8232           0.8141            0.7364
## Pos Pred Value            0.6601           0.6260            0.4542
## Neg Pred Value            0.8240           0.8141            0.7356
## Prevalence                0.3418           0.3320            0.3262
## Detection Rate            0.2260           0.2079            0.1478
## Detection Prevalence      0.3424           0.3321            0.3255
## Balanced Accuracy         0.7423           0.7201            0.5948
\end{verbatim}

Ahora ampliamos el alcance del sistema para que no solo diga si una casa
es ``cara o no'', sino que pueda clasificarla en tres niveles:
económica, intermedia o cara. Este modelo es más completo porque
entiende mejor la estructura de todo el mercado de Austin. Aunque es un
poco más complejo, nos da una visión mucho más detallada que ayuda a la
empresa a posicionar mejor cada propiedad frente a la competencia.

\section{Actividad 12: Comparación con métodos
anteriores}\label{actividad-12-comparaciuxf3n-con-muxe9todos-anteriores}

\begin{verbatim}
##                Metodo Exactitud
## 1 Regresion Logistica      0.72
## 2   Arbol de Decision      0.68
## 3       Random Forest      0.75
## 4         Naive Bayes      0.65
\end{verbatim}

Finalmente, ponemos frente a frente este nuevo modelo con los que
hicimos en semanas anteriores, como los árboles de decisión o los
bosques aleatorios. Observamos que la regresión logística es muy
competitiva: es mucho más rápida de procesar y más fácil de explicar que
un bosque aleatorio, manteniendo una precisión muy similar. Esto la
convierte en una opción excelente cuando se busca un sistema que sea
fácil de entender para los gerentes de la empresa sin sacrificar la
calidad de los resultados.

\end{document}
